\section{Limitations}

We list these as constraints on what the numbers above can be used for, not as caveats to soften them.

\paragraph{Scope of the models and policies.} Three released generation policies from two open model
families, all reading fixed 6-mers. Carbon's deployed logits-processor stack is unresolved, so the
policy a Carbon user obtains by default is characterised for the main revision we audited and not for
every configuration the release admits. Carbon-3B was not run.

\paragraph{Scope of the cohort.} The channel measurement uses all $24$ frozen public prompts, but every
detection and robustness result rests on $8$ prompts across four organisms. That is the resampling unit
for every interval we report, and it is why several intervals are coarse. Nothing here establishes
behaviour on other organisms, other genomic regions, or non-eukaryotic sequence.

\paragraph{Scope of the edit model.} Substitutions, insertions, deletions, crops, and reverse
complementation are applied by a synthetic protocol at declared rates. They are a stand-in for
sequencing and handling error, not a model of any particular platform's error profile, and the
adversarial attacks were never composed with them.

\paragraph{Biological proxies are proxies.} Every sequence statistic here is structural: composition,
complexity, or reading frame. None establishes function, viability, or safety, and we make no such
claim. In particular, a change in the longest open reading frame of a high-entropy sequence is evidence
that order changed, not that anything biological was lost, because such a frame arises largely by
chance.

\paragraph{The removal attack is unpriced.} We implemented both order-sensitive measures our design
named for this purpose and neither prices it. The composition proxies are blind by construction, the
independent-model score shifts by at most $\independentModelShift{}$, and the reading-frame measure
shows no directionally consistent change under any bounded rearrangement. We therefore report that a
6-mer permutation removes the watermark and that we cannot say what it costs. A reader should not read
the small measured shifts as evidence that the attack is cheap.

\paragraph{Statistical resolution.} Achieved false-positive rates are multiples of one over the number
of null trials, and the intervals cannot resolve below that. Several attack and many-query cells carry
eight trials, so they support statements about whether an effect is present, not precise rates. Where a
target rate is below the attainable resolution we say so rather than reporting it.

\paragraph{Exact marginals are not security.} One-step distribution preservation is proved and tested,
and it implies nothing about key reuse, adaptive queries, or spoofing. Our own key-reuse result is the
sharpest available demonstration: a construction with exactly correct marginals admits a complete
spoofing attack. Any cryptographic claim about this construction would need its own assumptions and
proof, and we make none.

\paragraph{Feasibility, not benchmarking.} Runtimes are single observations on one laptop under
uncontrolled load and on mains power; the same work runs about an order of magnitude slower on battery.
No peak-memory figure is reported, because the counters available to us cannot support a residency
claim.

\paragraph{One unresolved result.} On the base-marginal policy, two of nine held-out unkeyed
distinguisher tests reject nominally without clearing a Bonferroni correction across the family. We
report it as unresolved rather than as either a detection or a clean bill of health; resolving it needs
more matched draws than we have collected.
